% Generated draft: requires mathematical and source review.
\documentclass[11pt]{article}
\usepackage[margin=25mm]{geometry}
\usepackage{fontspec}
\setmainfont{Latin Modern Roman}
\usepackage{amsmath,amssymb,mathtools,unicode-math}
\setmathfont{Latin Modern Math}
\usepackage{enumitem,needspace,xurl,hyperref}
\setlength{\parindent}{0pt}
\setlength{\parskip}{5pt}
\setlength{\emergencystretch}{3em}
\newcommand{\N}{\mathbb{N}}
\newcommand{\Z}{\mathbb{Z}}
\newcommand{\Q}{\mathbb{Q}}
\newcommand{\R}{\mathbb{R}}
\newcommand{\C}{\mathbb{C}}
\newcommand{\F}{\mathbb{F}}
\newcommand{\A}{\mathbb{A}}
\newcommand{\PP}{\mathbb{P}}
\newcommand{\E}{\mathbb{E}}
\newcommand{\eps}{\varepsilon}
\newcommand{\dd}{\,\mathrm{d}}
\begin{document}
\section*{UnsolvedMath ID: 700797}
\textbf{Source problem code:} catalogue entry
\subsection{Definitions and mathematical statement}
The remaining question is part (b) of Kourovka Notebook Problem 11.105, attributed to G. E. Wall (1990). The following formulation uses the standard lower central series interpretation of the source's phrase ``associated Lie ring'' and the natural presentation induced by the free generators. These conventions are not explicitly specified in the supplied record and require review.

Let \(q=p^k\geq 4\), where \(p\) is prime and \(k\geq 1\) is an integer. Let \(\mathfrak B_q\) denote the Burnside variety consisting of all groups satisfying \(g^q=1\) for every element \(g\); the identity does not require each element to have order exactly \(q\). Let \(X\) be a free generating set, let \(F=F(X)\) be the absolutely free group on \(X\), and put
\[
N_q=\langle f^q\mid f\in F\rangle,\qquad G_q=F/N_q.
\]
Thus \(G_q\) is the relatively free group in \(\mathfrak B_q\) on \(X\), and \(N_q\) is fully invariant in \(F\), meaning that every group endomorphism of \(F\) maps \(N_q\) into itself.

For a group \(G\), define its lower central series by
\[
\gamma_1(G)=G,\qquad \gamma_{n+1}(G)=[\gamma_n(G),G]\quad(n\geq 1),
\]
and define the associated Lie ring by
\[
\mathcal L(G)=\bigoplus_{n\geq 1}\gamma_n(G)/\gamma_{n+1}(G),
\]
with addition from the abelian quotient groups and Lie bracket induced by group commutators. Let \(L=L_{\Z}(X)\) be the free Lie ring over \(\Z\) on \(X\). Sending each \(x\in X\) to its degree-one class in \(\mathcal L(G_q)\) gives a surjective Lie ring homomorphism
\[
\pi_q:L\longrightarrow\mathcal L(G_q).
\]
Set \(J_q=\ker\pi_q\), so that \(\mathcal L(G_q)\cong L/J_q\). An ideal \(J\) of \(L\) is fully invariant if
\[
\theta(J)\subseteq J\qquad\text{for every Lie ring endomorphism }\theta:L\to L.
\]
Here endomorphisms are not required to preserve the grading.

\textbf{Question (b).} For every prime power \(q\geq 4\) and every rank under consideration, is \(J_q\) fully invariant in \(L\)? Equivalently, does
\[
\theta(J_q)\subseteq J_q
\qquad\text{hold for every }\theta\in\operatorname{End}_{\mathrm{Lie}}(L)?
\]
The source says ``of given rank'' without specifying whether ranks must be finite or imposing a lower bound. The quantifier over ranks above is therefore provisional; no rank restriction is supplied.

\textbf{Archived part (a) and supplied qualification.} The original general question asks the same full invariance question for an arbitrary variety \(\mathfrak V\) of groups: if \(F/N\) is relatively free in \(\mathfrak V\), with \(N\) fully invariant in \(F\), and \(\mathcal L(F/N)\cong L/J\) is the corresponding natural presentation, must \(J\) be fully invariant in \(L\)? The supplied archive reports a negative answer for
\[
F/(F^2)^4,\qquad\mathfrak V=\mathfrak B_4\mathfrak B_2,
\]
attributed to D. Groves. In the usual power-subgroup notation, \(F^2=\langle f^2\mid f\in F\rangle\) and \((F^2)^4=\langle h^4\mid h\in F^2\rangle\). The record gives no rank for this example. This is a supplied report concerning part (a), not an independently verified result or a resolution of part (b).
\subsection{Short English statement}
For a relatively free group in the Burnside variety satisfying \(g^q=1\), where \(q\geq 4\) is a prime power, ask whether the kernel of the natural map from the free Lie ring to its lower central series Lie ring is preserved by every Lie ring endomorphism. The source reports that the analogous question for arbitrary group varieties has a negative answer, while retaining the Burnside question as the remaining part. The associated Lie ring convention, natural presentation, and intended range of ranks require confirmation.
\subsection{Sources}
\begin{itemize}
\item Supplied local UnsolvedMath registry record, permanent ID \(700797\), exact source problem code \texttt{catalogue entry}, title ``Kourovka Notebook Problem 11.105.'' Supplied upstream revision: \texttt{d371ee917ac9647fbb06cd6fc529deec6f3dda82}.
\item \emph{The Kourovka Notebook (Unsolved Problems in Group Theory)}, 11th issue (1990), Problem 11.105, proposed by G. E. Wall; supplied record describes an update in the 21st issue (2026), \texttt{arXiv:1401.0300v48}. The supplied collection description names E. I. Khukhro and V. D. Mazurov as editors. These bibliographic details are supplied attribution; the referenced text has not been independently consulted here.
\item Supplied Archive of Solved Problems entry for part (a): D. Groves, \emph{J. Algebra} \textbf{211}, no. 1 (1999), 15--25. The entry reports that \(J\) is not fully invariant for \(F/(F^2)^4\), corresponding to \(\mathfrak B_4\mathfrak B_2\). No article title or direct link is supplied.
\item Supplied status review dated \texttt{2026-10-16}: canonical label \texttt{partially\_solved}, evidence field ``published''; its note places solved part (a) in the archive and the remaining part in the main section. This is supplied dated curation, not an independently verified mathematical or literature assessment.
\end{itemize}
\end{document}
